\documentclass[11pt, a4paper]{article}

\usepackage[utf8]{inputenc}
\usepackage[T1]{fontenc}
\usepackage{amsmath, amssymb}
\usepackage{geometry}
\usepackage{booktabs}
\usepackage{array}
\usepackage{colortbl}
\usepackage{xcolor}
\usepackage{titlesec}
\usepackage{enumitem}
\usepackage{fancyhdr}
\usepackage{mdframed}

\geometry{top=2.2cm, bottom=2.2cm, left=2.5cm, right=2.5cm}

\definecolor{ucblue}{HTML}{1a237e}
\definecolor{ucbluelight}{HTML}{3949ab}
\definecolor{rowgray}{HTML}{e8eaf6}
\definecolor{codebg}{HTML}{f3f4f6}
\definecolor{accentred}{HTML}{c62828}
\definecolor{accentgreen}{HTML}{2e7d32}

\titleformat{\section}{\large\bfseries\color{ucblue}}{}{0em}{}[\vspace{-4pt}\rule{\linewidth}{0.8pt}]
\titlespacing{\section}{0pt}{12pt}{5pt}

\titleformat{\subsection}{\normalsize\bfseries\color{ucbluelight}}{}{0em}{}
\titlespacing{\subsection}{0pt}{8pt}{3pt}

\newmdenv[
  linecolor=ucbluelight, linewidth=1.5pt,
  topline=false, rightline=false, bottomline=false,
  leftmargin=8pt, innerleftmargin=10pt,
  innertopmargin=5pt, innerbottommargin=5pt,
  backgroundcolor=rowgray
]{nota}

\pagestyle{fancy}
\fancyhf{}
\lhead{\small\color{ucblue}\textbf{ICS1113} Optimización · Informe 3}
\rhead{\small\color{ucblue}Grupo 11}
\cfoot{\small\thepage}
\renewcommand{\headrulewidth}{0.4pt}

\setlength{\parindent}{0pt}
\setlength{\parskip}{4pt}

\begin{document}

\begin{center}
  {\LARGE\bfseries\color{ucblue} Resultados del modelo con datos reales}\\[3pt]
  {\normalsize Horizonte $T = 365$ días \quad Vida útil $L = 42$ días}
\end{center}
\vspace{-2pt}
{\color{ucblue}\rule{\linewidth}{1.5pt}}

\section{Resumen ejecutivo}

El modelo resolvió a óptimo en \textbf{35 minutos} (5{,}4M variables, 1{,}2M restricciones). La red satisface el 100\% de la demanda anual; el desperdicio resultante (7{,}8\%) se debe a sobre-oferta estructural de donaciones, no a fallas del modelo.

\begin{center}
\begin{tabular}{lr}
\toprule
\rowcolor{rowgray}\textbf{Indicador} & \textbf{Valor}\\
\midrule
Valor objetivo & 518.180,81 \\
Demanda total año & 52.659 bolsas \\
Demanda satisfecha & \textbf{\color{accentgreen} 52.659 (100\%)} \\
Demanda insatisfecha & 0 (0\%) \\
Bolsas vencidas & \textbf{\color{accentred} 4.577 (7,8\% de la oferta)} \\
Tiempo de resolución & 35 min \\
\bottomrule
\end{tabular}
\end{center}

\section{Desglose del objetivo}

\begin{center}
\begin{tabular}{lrr}
\toprule
\rowcolor{rowgray}\textbf{Término} & \textbf{Valor} & \textbf{\% del total}\\
\midrule
1. Calidad de uso $\sum[(L-u)+\tau_h+\alpha_g]\,y$ & 152.020,81 & 29,3\% \\
2. Demanda insatisfecha $p\sum s$ & 0,00 & 0,0\% \\
3. Vencimientos $\eta \sum w$ & \textbf{366.160,00} & \textbf{70,7\%} \\
\midrule
\textbf{Total} & \textbf{518.180,81} & 100\% \\
\bottomrule
\end{tabular}
\end{center}

El término dominante es \textbf{vencimientos}, no insatisfacción. El sistema no tiene problema de cobertura sino de gestión del excedente.

\section{Balance oferta--demanda}

\begin{center}
\begin{tabular}{lr}
\toprule
\rowcolor{rowgray}\textbf{Concepto} & \textbf{Bolsas}\\
\midrule
Oferta total & 58.564 \\
\quad Donaciones año & 57.818 \\
\quad Inventario inicial & 746 \\
Demanda total año & 52.659 \\
\midrule
\textbf{Exceso de oferta} & \textbf{5.905 \,(10,1\%)} \\
\quad Vencidas en el modelo & 4.577 \,(77,5\% del exceso) \\
\quad Remanente al día 365 & 1.328 \\
\bottomrule
\end{tabular}
\end{center}

Las donaciones programadas exceden la demanda en $\sim$10\% anual. El modelo no puede crear demanda: el excedente o se queda como inventario final o vence.

\section{Hallazgo crítico: pico de vencimientos al día 43}

\begin{nota}
\textbf{Día 43:} vencen 761 bolsas (17\% del total anual en un único día).
\end{nota}

\textbf{Causa.} El inventario inicial (746 bolsas) entra todo a $u = L = 42$ el día 1 por el supuesto 10 del PDF. Cumple su ciclo de envejecimiento y vence el día $L+1 = 43$.

\textbf{Por qué no se usa antes.} Las donaciones diarias bastan para cubrir la demanda diaria; el modelo prefiere usar bolsas frescas (término~1) y el inventario inicial envejece sin rotarse.

\textbf{Implicancia.} En la realidad, el inventario inicial de un banco está escalonado en distintas edades, no concentrado a $u=L$. \textit{Conviene mencionar el supuesto 10 como limitación en el informe.}

\section{Patrón temporal y por tipo}

Después del día 43, los vencimientos aparecen en \textbf{olas de 5--7 días} durante todo el año (días 50, 57, 78, 87, 106, 113, 134, 155, \ldots con 50--80 bolsas cada uno). Refleja el ciclo de envejecimiento del exceso crónico.

\begin{center}
\begin{tabular}{lrrr}
\toprule
\rowcolor{rowgray}\textbf{Tipo} & \textbf{Vencidas} & \textbf{Demanda año} & \textbf{Vencidas / Demanda}\\
\midrule
B+ & 1.503 & 4.445 & 33,8\% \\
AB+ & 617 & 972 & 63,5\% \\
O+ & 1.172 & 30.529 & 3,8\% \\
A+ & 1.026 & 13.859 & 7,4\% \\
O-- & 98 & 1.778 & 5,5\% \\
A-- & 85 & 755 & 11,3\% \\
B-- & 42 & 259 & 16,2\% \\
AB-- & 34 & 62 & 54,8\% \\
\bottomrule
\end{tabular}
\end{center}

\textbf{B+ y AB+} son los más castigados. AB+ tiene compatibilidad mínima (solo satisface a AB+, único receptor). B+ tiene relación donaciones/demanda muy desfavorable.

\section{Calidad del uso}

\begin{itemize}[leftmargin=*, itemsep=2pt]
  \item Costo medio por bolsa usada: $152.020,81 / 52.659 = \mathbf{2{,}89}$
  \item Composición típica: $(L-u)\approx 0$ + $\tau \approx 0{,}03$ + $\alpha \approx 2{,}8$
\end{itemize}

El modelo usa bolsas \textbf{muy frescas en promedio}: casi todas se utilizan el mismo día o al día siguiente de la donación.

\section{Validación contra literatura}

La tasa de desperdicio del \textbf{7,8\%} es coherente con la cita de Alsughayyir y cols. (2026, PLOS ONE) que reporta hasta \textbf{14,7\%} en hospitales sin optimización. El modelo reduce el desperdicio prácticamente a la mitad respecto a ese benchmark.

\section{Conclusiones para el informe}

\begin{enumerate}[leftmargin=*, itemsep=3pt]
  \item El modelo funciona correctamente: cubre el 100\% de la demanda sin holguras artificiales.
  \item El cuello de botella no es oferta sino \textbf{gestión del excedente} (donaciones $\sim$10\% mayores a la demanda).
  \item El supuesto 10 (inventario inicial todo a $u=L$) genera un pico artificial de vencimientos al día 43. Mencionar como limitación.
  \item La tasa de 7,8\% de desperdicio \textbf{mejora significativamente el 14,7\% de literatura}, validando el aporte del modelo.
  \item Recomendación al banco real: revisar la planificación de donaciones; un ajuste de $-5\%$ podría eliminar la mayoría de vencimientos sin comprometer cobertura.
\end{enumerate}

\section{Análisis de sensibilidad sugerido}

Para reforzar el informe, podrían correrse los siguientes escenarios:

\begin{itemize}[leftmargin=*, itemsep=2pt]
  \item \textbf{Donaciones $\times\,0{,}95$:} ¿se eliminan vencimientos manteniendo cobertura?
  \item \textbf{$\eta = 200$:} ¿cambia algo o el modelo ya está en su mínimo posible?
  \item \textbf{Inventario inicial escalonado en edades $1..L$:} ¿desaparece el pico del día 43?
\end{itemize}

\vfill
{\color{ucblue}\rule{\linewidth}{0.6pt}}
\begin{center}
  \small Modelo resuelto con Gurobi 10.0.3 · 1.168.458 restricciones, 5.445.800 variables · óptimo global alcanzado
\end{center}

\end{document}
